% BHU PYQ -- Manually transcribed & error-corrected
% Paper: MTB-606 : Number Theory
% Exam: B.A./B.Sc. (Hons) Semester VI Examination, 2013-14

\documentclass[11pt,a4paper]{article}
\usepackage[a4paper,top=2.5cm,bottom=2.5cm,left=2.8cm,right=2.8cm,headheight=14pt]{geometry}
\usepackage{amsmath,amssymb}
\usepackage[T1]{fontenc}\usepackage[utf8]{inputenc}
\usepackage{lmodern,microtype,fancyhdr,enumitem,hyperref,lastpage,parskip}

\hypersetup{colorlinks=true,urlcolor=black,linkcolor=black,
  pdftitle={MTB-606: Number Theory -- BHU B.A./B.Sc. Sem VI 2013-14}}
\pagestyle{fancy}\fancyhf{}
\renewcommand{\headrulewidth}{0.4pt}\renewcommand{\footrulewidth}{0.4pt}
\fancyhead[L]{\small\textit{Banaras Hindu University}}
\fancyhead[R]{\small\textit{MTB-606: Number Theory}}
\fancyfoot[C]{\small Page \thepage\ of \pageref{LastPage}}

\newlist{parts}{enumerate}{2}
\setlist[parts,1]{label=(\alph*),leftmargin=*,itemsep=5pt,topsep=3pt}
\setlist[parts,2]{label=(\roman*),leftmargin=*,itemsep=3pt,topsep=2pt}
\newcommand{\pts}[1]{\hfill\textit{[#1]}}

\begin{document}
\begin{center}
  {\large\bfseries BANARAS HINDU UNIVERSITY}\\[2pt]
  {\normalsize B.A./B.Sc. (Hons) --- Semester VI Examination, 2013-14}\\[6pt]
  \rule{0.8\linewidth}{0.5pt}\\[6pt]
  {\large\bfseries Paper No. MTB-606: Number Theory}\\[4pt]
  {\normalsize Time: 3 Hours \hfill Full Marks: 70}
\end{center}
\rule{\linewidth}{0.4pt}
\smallskip
\noindent\textit{Instructions: Answer \textbf{five} questions including Question~1 which is compulsory. Figures in right-hand margin indicate marks.}
\bigskip

\noindent\textbf{Question 1.}
\begin{parts}
  \item Show that for any positive integer $m$: $(ma, mb) = m(a,b)$. \pts{2}
  \item Find integers $x, y$ such that $9 = 657x + 963y$. \pts{2}
  \item Prove that if $n$ is composite, then it has a prime divisor $p$ satisfying $p \le \sqrt{n}$. \pts{2}
  \item If $2^n - 1$ is prime, prove that $n$ is prime. \pts{2}
  \item Show that $ax \equiv ay \pmod{m}$ if and only if $x \equiv y \pmod{\dfrac{m}{(a,m)}}$. \pts{2}
  \item If $n$ is composite and $n > 4$, prove that $(n-1)! \equiv 0 \pmod{n}$. \pts{2}
  \item Find all Pythagorean triples whose terms form an arithmetic progression. \pts{2}
\end{parts}

\medskip\rule{\linewidth}{0.2pt}

\medskip\noindent\textbf{Question 2.}
\begin{parts}
  \item If $g = \gcd(a,b)$, prove that there exist integers $x$ and $y$ such that $g = ax + by$. \pts{7}
  \item If $m > 0$, prove that $[ma, mb] = m[a,b]$. Also prove that $[a,b](a,b) = |ab|$. \pts{7}
\end{parts}

\medskip\rule{\linewidth}{0.2pt}

\medskip\noindent\textbf{Question 3.}
\begin{parts}
  \item Prove that if $(a,m) = (b,m) = 1$, then $(ab,m) = 1$. \pts{6}
  \item State and prove the Fundamental Theorem of Arithmetic. \pts{8}
\end{parts}

\medskip\rule{\linewidth}{0.2pt}

\medskip\noindent\textbf{Question 4.}
\begin{parts}
  \item State and prove the Chinese Remainder Theorem. \pts{5}
  \item Find all integers that have remainders $1$ or $2$ when divided by each of $3$, $4$, and $5$. \pts{9}
\end{parts}

\medskip\rule{\linewidth}{0.2pt}

\medskip\noindent\textbf{Question 5.}
\begin{parts}
  \item State and prove Wilson's theorem. \pts{7}
  \item Prove that if $(m,n) = 1$, then $\phi(mn) = \phi(m)\phi(n)$. Show that the condition $(m,n) = 1$ is necessary. \pts{7}
\end{parts}

\medskip\rule{\linewidth}{0.2pt}

\medskip\noindent\textbf{Question 6.}
\begin{parts}
  \item If $p$ is prime, prove that $x^2 \equiv -1 \pmod{p}$ has a solution if and only if $p = 2$ or $p \equiv 1 \pmod{4}$. \pts{7}
  \item Show that for $n \ge 1$: $\displaystyle\sum_{d \mid n} \phi(d) = n$. \pts{7}
\end{parts}

\medskip\rule{\linewidth}{0.2pt}

\medskip\noindent\textbf{Question 7.}
\begin{parts}
  \item Let $p$ be an odd prime and let $1 + \dfrac{1}{2} + \dfrac{1}{3} + \cdots + \dfrac{1}{p-1} = \dfrac{a}{b}$ where $\gcd(a,b) = 1$. Prove that $p \mid a$. \pts{7}
  \item Find all integral solutions of $32x + 6y + 15z = 1$. \pts{7}
\end{parts}

\vfill\rule{\linewidth}{0.4pt}
\begin{center}\small
  \href{https://bhu-exam.vercel.app/}{Click here to visit our website}\\[4pt]
  \href{https://me-aryan.vercel.app/}{Click here to visit our contributor}\\[4pt]
  \href{https://quashan-me.vercel.app/}{Click here to visit our contributor}
\end{center}
\end{document}
